% ==STYLE-BEGIN==
\documentclass[UTF8,11pt]{ctexart}
\usepackage{amsmath,amssymb,amsthm}
\usepackage[a4paper,margin=2.4cm,top=2.3cm,bottom=2.3cm]{geometry}
\usepackage{fancyhdr}
\usepackage{enumitem}
\usepackage{hyperref}
\hypersetup{hidelinks}

% ---- 版式 ----
\linespread{1.15}
\setlength{\parskip}{0.35em}
\setlength{\parindent}{1.6em}
\setlist{nosep,leftmargin=2em}

% ---- 页眉 ----
\pagestyle{fancy}
\fancyhf{}
\fancyhead[L]{\small\coursename}
\fancyhead[R]{\small\thepage}
\renewcommand{\headrulewidth}{0.4pt}

% ---- 定理环境：中文用正体，标题加粗，按小节共用一条流水号 ----
\newtheoremstyle{cnthm}{0.7em}{0.7em}{\normalfont}{}{\bfseries}{}{0.5em}{}
\theoremstyle{cnthm}
\newtheorem{thm}{定理}[section]
\newtheorem{lem}[thm]{引理}
\newtheorem{prop}[thm]{命题}
\newtheorem{cor}[thm]{推论}
\newtheorem{defn}[thm]{定义}
\newtheorem{rem}[thm]{注}
\newtheorem{exa}[thm]{例}
\newtheorem{conj}[thm]{猜想}
\renewcommand{\proofname}{\bfseries 证明}
\providecommand{\coursename}{}
% ==STYLE-END==
\renewcommand{\coursename}{解析几何\,·\,0917}
\title{习题课 0917}
\date{}
\begin{document}
\maketitle

\section{\ 有轴平面束}

\[
\pi_1:\ A_1x+B_1y+C_1z+D_1=0,\qquad
\pi_2:\ A_2x+B_2y+C_2z+D_2=0 .
\]
$\pi_1$ 与 $\pi_2$ 的交线为 $l$，设 $\pi$ 是任意一个过直线
$l$ 的平面，取 $P(x_0,y_0,z_0)\in\pi-l$，此时有
\[
A_1x_0+B_1y_0+C_1z_0+D_1\neq0
\quad\text{or}\quad
A_2x_0+B_2y_0+C_2z_0+D_2\neq0 .
\]
不妨令
\[
\lambda=A_2x_0+B_2y_0+C_2z_0+D_2,\qquad
\mu=-(A_1x_0+B_1y_0+C_1z_0+D_1),
\]
则方程
\[
\lambda(\pi_1)+\mu(\pi_2)=0
\]
表示一个平面，且该平面经过 $P$ 和 $l$，说明该平面即为
$\pi$，于是 $\pi$ 的方程为上式。

$\lambda$ 与 $\mu$ 至少有一个不为 $0$。

\section{\ 化为点向方程：以 (4) 为例}

\[
\frac{x}{1}=\frac{y-1}{0}=\frac{z+2}{0},
\]
某一个分母为 $0$ 应理解为它的分子为 $0$。
点向方程的分母可以为 $0$。

\end{document}

\end{document}
